\begin{longtable}{TD}
\caption{Profiling and observation-policy terms.}
\label{tab:terms-profiling-observation-policy}\\

\toprule
Term & Definition \\
\midrule
\endfirsthead

\toprule
Term & Definition \\
\midrule
\endhead

\bottomrule
\endfoot

\termrow{profiling}{
The process of collecting, interpreting, and using observations to estimate
task behavior. In this paper, profiling is broader than CPU profiling.
}

\termrow{task profiling}{
Profiling logical task attempts in queue-worker systems in order to estimate
resource demand, worker occupancy, uncertainty, and operational risk.
}

\termrow{observation strategy}{
A decision about what to observe, where to observe it, how often to observe it,
and under what overhead budget.
}

\termrow{profiling policy}{
The selected observation and instrumentation strategy for a task kind, task
attempt, or execution context.
}

\termrow{accounting}{
Recording basic execution facts, such as start time, finish time, status,
duration, input size, output size, resource usage, or occupancy.
}

\termrow{minimal accounting}{
A low-overhead accounting level that records only essential task facts, such as
start time, finish time, status, and basic size information.
}

\termrow{detailed accounting}{
An accounting level that records more detailed resource, occupancy, stage,
dependency, and outcome information.
}

\termrow{tracing}{
Capturing execution structure and timing through traces and spans. Tracing is
an observation method, not a behavior model by itself.
}

\termrow{sampled tracing}{
Tracing only a subset of tasks, attempts, or traces in order to control
overhead.
}

\termrow{tail-aware profiling}{
A profiling strategy that preferentially observes slow, failed, high-risk, or
tail-heavy attempts.
}

\termrow{diagnostic profiling}{
Deep and usually expensive profiling used to investigate unexplained behavior,
high residual error, drift, or model failure.
}

\termrow{isolated execution profiling}{
Profiling in a more controlled or isolated execution environment in order to
reduce interference or better attribute behavior.
}

\termrow{cooperative instrumentation}{
Instrumentation in which worker or task code explicitly emits structured
observations such as spans, timers, byte counts, record counts, dependency
calls, stage boundaries, and domain-specific signals.
}

\termrow{sampling}{
Selecting a subset of tasks, attempts, events, spans, traces, or profile samples
for observation or retention.
}

\termrow{head sampling}{
A trace-sampling strategy that decides whether to sample before the full trace
is observed.
}

\termrow{tail sampling}{
A trace-sampling strategy that decides whether to sample after seeing all or
most of a trace.
}

\termrow{observation point}{
A point in the task lifecycle where observations can be collected, such as
enqueue, dequeue, reserve, start, dependency call, retry, commit, acknowledge,
or finish.
}

\termrow{profiling budget}{
The allowed overhead for profiling and telemetry, including CPU, memory,
allocation, storage, network, and cardinality costs.
}

\termrow{hot path}{
The latency- or throughput-sensitive path of task execution where added
instrumentation overhead can affect production behavior.
}

\end{longtable}
